\documentclass[sigconf,nonacm]{acmart}

\settopmatter{printfolios=true}
\setlength{\emergencystretch}{1em}

\begin{document}

\title{One GPU, 0.1 Trillion Tokens: Pretraining a Sub-Billion-Parameter Code Model}

\author{Indrajeet Aditya Roy}
\email{roy.indrajeet009@gmail.com}
\affiliation{%
  \institution{}
  \country{}}

\begin{abstract}
Code large language models (LLMs) with fewer than one billion parameters are typically pretrained on trillions of
tokens. This paper presents a generative LLM for Python and C++ code generation: a GPT-style decoder-only transformer
pretrained from scratch on a single GPU. The model has 522M parameters and is trained on 103B tokens of Python and C++
code, web text and math, first with a context length of 8,192 tokens at the file level and then with 32,768 tokens at
the repository level. Pretraining takes 76.3 hours on one NVIDIA GH200 GPU.
\end{abstract}

\keywords{code language models, pretraining, small language models, training efficiency, code generation}

\maketitle

\section{Introduction}

Large language models (LLMs) are now strong at code generation and software engineering tasks. For example, Claude
Opus 5.5 resolves 89.9\% of the tasks in SWE-bench Pro~\cite{anthropic2026opus55}, a benchmark of tasks from actively
maintained repositories that require large, multi-file changes~\cite{deng2025swebenchpro}. Other frontier LLMs
released in 2026 include OpenAI's GPT-6 Astra~\cite{openai2026gpt6astra} and DeepSeek-V4~\cite{deepseek2026v4}.

These capabilities come from scale. An LLM is pretrained as an autoregressive language model: it learns to predict
the next token of its training data by minimizing the cross-entropy training loss. Scaling laws show that this loss
decreases as a power law in the number of model parameters, the number of training tokens and the training
compute~\cite{kaplan2020scaling}. For a fixed compute budget, the training loss is lowest when model size and training
tokens are scaled equally~\cite{hoffmann2022chinchilla}: the compute-optimal Chinchilla model has 70B parameters and is
trained on 1.4T tokens, about 20 tokens per parameter. Frontier LLMs are therefore trained on large GPU clusters.
DeepSeek-V3, a mixture-of-experts (MoE) LLM with 671B total and 37B activated parameters, was pretrained on 14.8T
tokens on 2,048 H800 GPUs and used 2.788M GPU hours for its full training~\cite{deepseek2024v3}. DeepSeek-V4-Pro has
1.6T total and 49B activated parameters and was pretrained on more than 32T tokens~\cite{deepseek2026v4}. The system
cards of the Claude and GPT models do not report model size, training tokens or training
compute~\cite{anthropic2026opus55,openai2026gpt6astra}.

Code LLMs follow the same trend. DeepSeek-Coder was pretrained from scratch on 2T tokens~\cite{guo2024deepseekcoder},
StarCoder2 on 3.3--4.3T tokens~\cite{lozhkov2024starcoder2} and Seed-Coder on 6T tokens~\cite{seed2025coder}.
DeepSeek-Coder-V2 and Qwen2.5-Coder continue pretraining from general-purpose LLMs on an additional 6T and 5.5T tokens
of code-focused data~\cite{deepseek2024coderv2,hui2024qwen25coder}. Their instruction-tuned and reasoning models are
then post-trained with supervised fine-tuning (SFT), followed by preference optimization such as direct preference
optimization (DPO)~\cite{rafailov2023dpo,hui2024qwen25coder,seed2025coder} or reinforcement learning
(RL)~\cite{deepseek2024coderv2,seed2025coder}. Small code LLMs are trained far beyond the compute-optimal number of
tokens: StarCoder2-3B is trained on 3.1T tokens, about 1,000 tokens per parameter~\cite{lozhkov2024starcoder2}, and all
Qwen2.5-Coder models, including the 0.5B model, on over 5.5T tokens~\cite{hui2024qwen25coder}. All Qwen3 models,
including the 0.6B model, are pretrained on 36T tokens~\cite{yang2025qwen3}.

Training budgets of this size are not feasible on a single GPU. The 522M-parameter model in this work trains at 0.38M
tokens per second on one GPU; at this training throughput, the 5.5T tokens of Qwen2.5-Coder would take more than five
months. This paper studies whether a code LLM pretrained and post-trained on a single GPU, with far fewer tokens, can
reach code generation performance comparable to the smallest existing code LLMs. Code generation is measured by
functional correctness: pass@1 is the fraction of problems whose generated program passes all unit tests, with greedy
decoding or averaged over sampled programs~\cite{chen2021codex}.

Scaling laws predict a higher training loss for a model trained on fewer tokens~\cite{kaplan2020scaling}. At small
scale, however, the focus and quality of the training data can substitute for its quantity: phi-1, a 1.3B-parameter
code model pretrained on about 50B tokens drawn from 7B tokens of filtered and synthetic data, reaches 29\% pass@1 on
HumanEval, and 50.6\% after fine-tuning on fewer than 200M tokens of exercises~\cite{gunasekar2023textbooks}. This
paper therefore makes three predictions. First, a 522M-parameter model pretrained on 103B tokens concentrated on
Python and C++ reaches pass@1 comparable to the smallest existing code LLMs on function-level benchmarks: HumanEval,
MBPP and their EvalPlus versions~\cite{chen2021codex,austin2021mbpp,liu2023evalplus}. Second, in post-training, SFT
provides most of the gain and RL adds little: for sub-billion-parameter models fully fine-tuned on SFT data,
additional RL brings no significant improvement~\cite{zhao2025mobilellmr1}. Third, competitive programming remains in
the low single digits: after reasoning post-training, MobileLLM-R1-360M scores 5.1 and Qwen2.5-0.5B-Instruct 3.6
pass@1 on LiveCodeBench v6~\cite{zhao2025mobilellmr1,jain2024livecodebench}.

The paper makes three contributions:
\begin{itemize}
  \item A 522M-parameter, GPT-style decoder-only LLM for Python and C++ code generation, pretrained from scratch on
  one NVIDIA GH200 GPU: 103B tokens (528 tokens per non-embedding parameter) in 76.3 hours, from a corpus of code, web
  text and math curated from 18 sources in 9 Hugging Face datasets, with a file-level stage at a context length of
  8,192 tokens, a repository-level stage at 32,768 tokens and fused GPU kernels.
  \item Post-training on the same GPU: SFT on reasoning traces, which teaches the model to reason before it answers,
  followed by RL with rewards from unit tests.
  \item An evaluation of every stage with each benchmark's official tool, in Python and C++, against the published
  scores of sub-billion-parameter models, with 95\% confidence intervals.
\end{itemize}

\section{Data}
\label{sec:data}

The model is trained on three datasets, one per stage: a pretraining corpus of code, web text and mathematics,
reasoning traces for supervised fine-tuning, and programming problems with unit tests for reinforcement learning. This
section describes each in the order it is used: where its text comes from, how it is cleaned, decontaminated and
split, and, for pretraining, how its sources are mixed. Tokenization, packing and fill-in-the-middle, which turn this
text into training sequences, are described separately.

\subsection{Pretraining corpus}
\label{sec:data-pretrain}

\paragraph{Sources.} The corpus combines 18 sources from 9 Hugging Face datasets (Table~\ref{tab:sources}). Its code
is Python and C++ only, the two languages the model targets. Repository files come from Stack-Edu, StarCoder2's data
filtered for educational, well-documented code~\cite{benallal2025smollm2}, and from UltraData's classifier-selected
(L2) subsets of The Stack v3 and of its own code collection~\cite{tu2026ultradatacode}. Synthetic code adds the
exercises UltraData writes from its algorithmic L2 files (L3)~\cite{tu2026ultradatacode}, and Nemotron's Python files
rewritten by Qwen3-32B and Python files transpiled into C++~\cite{nvidia2025nemotron3nano}. Smaller sources contribute
functions from open-source projects (CodeXGLUE~\cite{lu2021codexglue}), Kaggle notebooks~\cite{lozhkov2024starcoder2}
and DeepSeek-R1 solutions to Codeforces problems (Codeforces-CoTs, finished solutions only). Natural language comes
from the web text of Nemotron-ClimbMix~\cite{diao2025climb} and the synthetic mathematics of UltraData-Math L3 in four
formats~\cite{wang2026ultradatamath}. The sources the mixture draws only in part (synthetic code, web text and
mathematics) are downloaded only as far as the mixture needs them; the others are used whole.

\begin{table*}
  \caption{Pretraining sources after cleaning and decontamination.}
  \label{tab:sources}
  \small
  \begin{tabular}{lllr}
    \toprule
    Group & Dataset and content & Language & Train documents \\
    \midrule
    Repository files & Stack-Edu: StarCoder2 data filtered for educational code & Python, C++ & 24.5M, 16.1M \\
     & UltraData-Stack-v3-L2: The Stack v3 files selected by classifiers & Python, C++ & 17.4M, 12.4M \\
     & UltraData-Code L2: code files selected by quality classifiers & Python, C++ & 42.0M, 13.1M \\
    Synthetic code & UltraData-Code L3: exercises written from algorithmic files & Python, C++ & 3.8M, 4.3M \\
     & Nemotron-Pretraining-Code-v2: files rewritten; Python transpiled into C++ & Python, C++ & 18.0M, 5.1M \\
    Functions & CodeXGLUE code-to-text: functions from open-source projects & Python & 0.25M \\
    Notebooks & Kaggle notebooks as text & Python & 0.57M \\
    Reasoning traces & Codeforces-CoTs: DeepSeek-R1 solutions to Codeforces problems & Python, C++ & 0.02M \\
    Web text & Nemotron-ClimbMix & -- & 15.6M \\
    Mathematics & UltraData-Math L3: synthetic mathematics in four formats & -- & 3.4M \\
    \midrule
    Total & & & 176.6M \\
    \bottomrule
  \end{tabular}
\end{table*}

\paragraph{Standardization and cleaning.} Every document is normalized as a compiler reads a source file: line
endings become LF and a leading byte-order mark is removed; the text is otherwise kept as published. Each Python file
of a code source must parse as Python~3.12. A file that does not is repaired by the smallest combination of three
fixes that makes it parse (The Stack v3's anonymization placeholders replaced by bare words, tab indentation expanded,
Python~2 syntax converted by 2to3) and dropped if none does. Exact duplicates are removed across all sources, keeping
the first copy in the order of Table~\ref{tab:sources}.

\paragraph{Decontamination.} A training document is removed when it shares a 10-gram of whitespace-separated words
with a benchmark text, the rule of Qwen2.5-Coder and DeepSeek-Coder~\cite{hui2024qwen25coder,guo2024deepseekcoder}.
The benchmark texts are the prompts and reference solutions of every task the evaluation scores, 1,428 in all:
HumanEval+ and MBPP+~\cite{liu2023evalplus}, HumanEval-X C++, the C++ prompts of MultiPL-E~\cite{cassano2023multiple}
and the HumanEval infilling programs~\cite{bavarian2022fim}. The filter removes 491,203 training documents (0.28\%),
matching every HumanEval+ task, 372 of the 378 MBPP+ tasks and 157 of the 164 HumanEval-X C++ references. It also
removes common code: 8,539 documents of one source share a 10-gram with the reference solution of a primality test
(HumanEval/31). Validation documents are not filtered, so likelihood measurements keep one fixed set.

\paragraph{Split and repository samples.} 1\% of the corpus is held out for validation by a hash of the unit a document
belongs to: its repository for repository files, so that no repository straddles the split, its problem for Codeforces
solutions, and the document itself otherwise. The files of each repository are grouped into samples connected by
imports and ordered so that every file follows the files it imports; the repository-level stage trains on these
samples. The curated corpus holds 176.6M training and 1.8M validation documents.

\paragraph{Mixture.} Pretraining runs in two stages with the mixtures of Table~\ref{tab:mixture}. Code makes up 87\% of
each stage, DeepSeek-Coder's share of source code~\cite{guo2024deepseekcoder}; its 13\% of natural language is taken
here as web text (10\%) and mathematics (3\%). The mixture keeps natural language although code alone scores higher on
code: in Qwen2.5-Coder's comparison, code only reached 49.8 and 40.3 on two code benchmarks against 48.3 and 38.3 for
70\% code, 20\% text and 10\% mathematics, which led on mathematics and general benchmarks~\cite{hui2024qwen25coder}.
Every problem the model solves is stated in natural language, and after post-training it reasons in natural language;
the mixture therefore keeps a natural-language share near DeepSeek-Coder's, rather than none. Synthetic code is 30\% of
the code, the share at which good mixtures of rephrased synthetic web text converge~\cite{kang2025synthetic}; that
study does not cover code, and the share is not ablated here. The notebooks, functions and Codeforces solutions are
read exactly once, each weighted by its size, since repeating a small part of the data many times degrades a
model~\cite{hernandez2022repeated}; within a group, sources are weighted by their sizes. The file-level stage trains on
single documents at a context of 8,192 tokens; the repository-level stage trains on repository samples at 32,768
tokens for the last 5.45\% of the steps, Qwen2.5-Coder's two contexts and its repository-level share (300B of 5.5T
tokens)~\cite{hui2024qwen25coder}. The repository-level stage keeps the shares of web text and mathematics, as
OpenCoder's annealing keeps most of its original data distribution to avoid forgetting~\cite{huang2024opencoder}.

\begin{table}
  \caption{The pretraining mixture: each group's share of a stage and its tokens.}
  \label{tab:mixture}
  \small
  \begin{tabular}{lrr}
    \toprule
    Group & Share & Tokens \\
    \midrule
    \multicolumn{3}{l}{\emph{File-level stage, context 8,192}} \\
    Repository files & 59.3\% & 57.8B \\
    Synthetic code & 26.1\% & 25.4B \\
    Notebooks & 1.6\% & 1.5B \\
    Functions & $<$0.1\% & 0.05B \\
    Web text & 10.0\% & 9.7B \\
    Mathematics & 3.0\% & 2.9B \\
    \midrule
    \multicolumn{3}{l}{\emph{Repository-level stage, context 32,768}} \\
    Repository samples & 81.5\% & 4.6B \\
    Codeforces solutions & 5.5\% & 0.3B \\
    Web text & 10.0\% & 0.6B \\
    Mathematics & 3.0\% & 0.2B \\
    \midrule
    Total & & 103.0B \\
    \bottomrule
  \end{tabular}
\end{table}

\subsection{Supervised fine-tuning data}
\label{sec:data-sft}

\paragraph{Sources.} SFT teaches the model to reason before it answers, on released reasoning traces
(Table~\ref{tab:sft}). Python traces at volume come from the code data of Nemotron-Cascade's two SFT stages, solutions
by DeepSeek-R1 and DeepSeek-R1-0528~\cite{wang2025nemotroncascade}. Harder competitive-programming problems with
standard input and output, and the only C++ traces, come from Nemotron's competitive-programming SFT set. Prompts in
the form of HumanEval and MBPP, a function to complete, with short reasoning come from that set's Exercism exercises and
from KodCode, whose DeepSeek-R1 solutions all pass their problem's tests~\cite{xu2025kodcode}. KodCode's questions
flagged as similar to benchmarks (cosine similarity above 0.95 to HumanEval, MBPP, BigCodeBench or LiveCodeBench) are
left out.

\paragraph{Filtering.} Each trace is split into its problem statement, its reasoning and its final answer. A trace is
dropped when it lacks reasoning or an answer, when its reasoning is under 20 tokens, or when the whole document exceeds
the 32,768-token context; it is dropped rather than cut, since a cut trace loses its answer, which removes about half of
the competitive-programming traces. The pretraining decontamination rule is applied to statement, reasoning and answer
against the same benchmark texts and MBPP's 500 original test problems; it removes every trace of 5.5\% and 5.2\% of
the problems of Nemotron-Cascade's two stages and 5.6\% of KodCode's rows. Exact duplicates are removed, and a problem
keeps at most two traces in the four long-trace groups, a cap set by the training time (two epochs within 12.5B
tokens); the Exercism and KodCode problems have about one trace each.

\paragraph{Split.} The validation split holds out 1\% of the problems rather than of the prompts. One problem recurs
across sources under different wrappers, such as a language header or a function header written in two forms, so a
problem is identified by its statement with these wrappers removed, lowercased and reduced to its words. In a build
without KodCode, a split by exact prompt left 928 of 2,921 validation problems with their statement in the training
split; a split by problem left none.
Outside KodCode, no SFT prompt shares more than 18.6\% of the distinctive 8-word sequences of any of the 454
LiveCodeBench problems the evaluation uses~\cite{jain2024livecodebench}.

\begin{table}
  \caption{SFT data after filtering: training documents (distinct prompts) and tokens. All traces are Python except
  the C++ competitive-programming ones.}
  \label{tab:sft}
  \small
  \begin{tabular}{lrr}
    \toprule
    Group & Documents (prompts) & Tokens \\
    \midrule
    Nemotron-Cascade 1 & 227,585 (149,076) & 1.55B \\
    Nemotron-Cascade 2 & 147,317 (74,005) & 1.72B \\
    Competitive, Python & 48,007 (24,540) & 0.55B \\
    Competitive, C++ & 48,017 (24,500) & 0.56B \\
    Exercism & 76,375 (75,341) & 0.08B \\
    KodCode & 250,649 (250,646) & 1.19B \\
    \midrule
    Total & 797,950 (539,127) & 5.65B \\
    \bottomrule
  \end{tabular}
\end{table}

\subsection{Reinforcement learning data}
\label{sec:data-rl}

RL needs problems whose tests decide correctness. Its problems are the KodCode rows of the SFT training split, 250,646
distinct prompts, each with KodCode's tests: a pytest file for a function to complete, or standard-input cases for a
program~\cite{xu2025kodcode}. A test rewards a correct answer only if it calls what the prompt asks for, so two kinds of
rows are left out: prompts that do not name every function their tests import, where a correct answer under another
name fails, and problems that test library APIs, whose library versions differ from KodCode's (45 of 171 of their
reference solutions fail here). This leaves 180,911 problems: 107,244 functions to complete, 52,740 functions described
in prose and 20,927 standard-input programs. The test harness passes 2,995 of 3,000 sampled reference solutions, and a
deliberately wrong program passes the tests of 3 of them; a pass on such a problem is not counted. RL is evaluated on
the KodCode rows of the SFT validation split, held out from the training problems by question and by test file, since
the SFT split separates prompts and one question can appear in several styles: 1,115 of 2,547 rows remain, 785 of them
usable, from which 512 problems are drawn, stratified by style and difficulty.

\bibliographystyle{ACM-Reference-Format}
\bibliography{references}

\end{document}
